\chapter*{Abstract}
\thispagestyle{empty}

Diabetes mellitus, with critical emphasis on Type 2, has established itself as one of the most pressing health emergencies in the global scenario, being a chronic disease intrinsically associated with modern lifestyle factors that imposes a severe and growing burden on public health systems. Given the vastness and heterogeneity of data generated in monitoring this pathology, traditional analytical approaches often prove limited, justifying the emergence of artificial neural networks as indispensable and promising technological tools for detection and prevention. Unlike linear statistical methods, these advanced algorithms possess the capability to process large volumes of clinical, demographic, and behavioral information identifying complex, hidden, and nonlinear patterns to classify individual disease risk with high precision. The strategic integration of these predictive models into clinical decision support systems represents a significant advancement, as it not only enhances the achievement of faster and more accurate diagnoses but also underpins precision medicine through the development of personalized early intervention strategies. Ultimately, the application of this computational intelligence aims to reverse current epidemiological trends, drastically reducing the incidence, long-term complications, and socioeconomic impact of diabetes in contemporary society.
\\\\
Keywords: Neural networks, diabetes, prediction, artificial intelligence, machine learning.